% GENERATED FILE. Edit canonical lessons, then run npm run generate:books.
% canonical-source-hash: fnv1a64:9c289776a419c54f
% canonical-lessons: HI-C279-khakar, HI-C279-akhir, HI-C279-pahlephir, HI-C279-chuttikadin, HI-R279-a-story-in-order

\chapter{Telling It in Order}
\label{ch:hi-telling-it-in-order}
\hlchaptermodality{279}

\begin{chapteropening}
\textbf{By the end of this chapter:} \emph{I can tell a story in order with pahle, phir, uske bād and ākhir meṁ, join actions with -kar, and read a short story.}

Complete HI-R279-a-story-in-order: tell a short story in order, after reading Mira's holiday in the village.
\end{chapteropening}

\section[\texorpdfstring{\hi{खाना} \hi{खाकर} (khānā khākar)}{खाना खाकर (khānā khākar)}]{\hi{खाना} \hi{खाकर} (khānā khākar) — one action, then the next}

\label{lesson:HI-C279-khakar}



\begin{quote}
You met \textbf{-\hi{कर}} with \textbf{\hi{देखकर}} and \textbf{\hi{करके}}: \textbf{\hi{मैं} \hi{किताब} \hi{देखकर} \hi{आता} \hi{हूँ}}, \textquotedblleft{}I'll look at the book and come\textquotedblright{}. Say it once. Then \textquotedblleft{}she came\textquotedblright{}. (\emph{vah āī}.)

This lesson takes \textbf{-\hi{कर}} into the past, where stories are told.
\end{quote}

\subsection*{You'll want to know: -\hi{कर} in a story}
\begin{quote}
\textbf{\hi{मैंने} \hi{हाथ} \hi{धोकर} \hi{खाना} \hi{खाया।}}
\end{quote}

\begin{quote}
\emph{maiṁne hāth dhokar khānā khāyā.} — \textquotedblleft{}I washed my hands and ate.\textquotedblright{}
\end{quote}

One doer, two actions, in the order they happened. The first drops its \textbf{-\hi{ना}} and takes \textbf{-\hi{कर}}; only the last verb carries the past. \textbf{\hi{करना}} still gives \textbf{\hi{करके}}.

The last verb decides \textbf{\hi{ने}}. If it takes a thing, the doer takes \textbf{\hi{ने}}: \textbf{\hi{मैंने} \hi{हाथ} \hi{धोकर} \hi{खाना} \hi{खाया}}. If it does not, there is no \textbf{\hi{ने}}, though something was eaten on the way: \textbf{\hi{मैं} \hi{खाना} \hi{खाकर} \hi{सो} \hi{गया}} (\emph{maiṁ khānā khākar so gayā}), \textquotedblleft{}I ate and fell asleep\textquotedblright{}.

\begin{itemize}
\raggedright
  \item \textbf{\hi{वह} \hi{सुबह} \hi{उठकर} \hi{नहाया।}} (\emph{vah subah uṭhkar nahāyā.}) — He got up in the morning and bathed.
  \item \textbf{\hi{मीरा} \hi{ने} \hi{घर} \hi{आकर} \hi{चाय} \hi{पी।}} (\emph{Mīrā ne ghar ākar cāy pī.}) — Mira came home and drank tea.
\end{itemize}

\subsection*{You'll want to know: after doing}
\begin{quote}
\textbf{\hi{नाश्ते} \hi{के} \hi{बाद} \hi{हम} \hi{बाज़ार} \hi{गए।}}
\end{quote}

\begin{quote}
\emph{nāśte ke bād ham bāzār gae.} — \textquotedblleft{}After breakfast we went to the market.\textquotedblright{}
\end{quote}

\textbf{\hi{के} \hi{बाद}}, \textquotedblleft{}after\textquotedblright{}, follows a noun in its oblique form: \textbf{\hi{नाश्ते} \hi{के} \hi{बाद}}. It follows a verb too, as \textbf{\hi{के} \hi{लिए}} does in \textbf{\hi{करने} \hi{के} \hi{लिए}}: the \textbf{-\hi{ना}} becomes \textbf{-\hi{ने}}.

\begin{itemize}
\raggedright
  \item \textbf{\hi{खाना} \hi{खाने} \hi{के} \hi{बाद}} (\emph{khānā khāne ke bād}) — after eating
  \item \textbf{\hi{काम} \hi{करने} \hi{के} \hi{बाद} \hi{मैं} \hi{घर} \hi{आया।}} (\emph{kām karne ke bād maiṁ ghar āyā.}) — After work I came home.
\end{itemize}

\textbf{-\hi{कर}} needs one doer for both actions. When the doers differ, tell it in two sentences, or with \textbf{\hi{जब} … \hi{तब} …}.

\subsection*{Your turn}
Say these aloud:

\begin{itemize}
\raggedright
  \item \textquotedblleft{}I washed my hands and ate\textquotedblright{} — \emph{maiṁne hāth dhokar khānā khāyā}
  \item \textquotedblleft{}Mira came home and drank tea\textquotedblright{} — \emph{Mīrā ne ghar ākar cāy pī}
\end{itemize}

\begin{tcolorbox}[breakable,colback=teal!4,colframe=teal!35!black,title={Before you move on}]
How many verbs in a \textbf{-\hi{कर}} sentence carry the past? (One: the last.) Which verb decides \textbf{\hi{ने}}? (The last: \textbf{\hi{मैंने} \hi{हाथ} \hi{धोकर} \hi{खाना} \hi{खाया}}, but \textbf{\hi{मैं} \hi{खाना} \hi{खाकर} \hi{सो} \hi{गया}}.) \textquotedblleft{}After breakfast\textquotedblright{}? (\textbf{\hi{नाश्ते} \hi{के} \hi{बाद}}.) \textquotedblleft{}After eating\textquotedblright{}? (\textbf{\hi{खाना} \hi{खाने} \hi{के} \hi{बाद}}.)
\end{tcolorbox}
\section[\texorpdfstring{\hi{आख़िर} (ākhir)}{आख़िर (ākhir)}]{\hi{आख़िर} (ākhir) — at last, in the end}

\label{lesson:HI-C279-akhir}



\begin{quote}
Before the new one: say \textquotedblleft{}I ate and fell asleep\textquotedblright{}. (\emph{maiṁ khānā khākar so gayā}.) A story needs a word for its last step, and for a wait that ends.
\end{quote}

\subsection*{You'll want to know: \hi{आख़िर}}
\textbf{\hi{आख़िर}} — \emph{ākhir} — \textquotedblleft{}at last, finally\textquotedblright{}.

\begin{quote}
\textbf{\hi{दो} \hi{घंटे} \hi{के} \hi{बाद} \hi{आख़िर} \hi{ट्रेन} \hi{आई।}}
\end{quote}

\begin{quote}
\emph{do ghaṇṭe ke bād ākhir ṭren āī.} — \textquotedblleft{}After two hours the train came at last.\textquotedblright{}
\end{quote}

With \textbf{\hi{में}} it closes a story: \textbf{\hi{आख़िर} \hi{में}} (\emph{ākhir meṁ}), \textquotedblleft{}in the end\textquotedblright{}.

\begin{quote}
\textbf{\hi{आख़िर} \hi{में} \hi{हम} \hi{घर} \hi{पहुँचे।}}
\end{quote}

\begin{quote}
\emph{ākhir meṁ ham ghar pahũce.} — \textquotedblleft{}In the end we reached home.\textquotedblright{}
\end{quote}

\begin{cousinweb}
\textbf{\hi{आख़िर}} (\emph{ākhir}) is usually traced to Arabic \emph{ākhir}, \textquotedblleft{}last, the end\textquotedblright{}, which reached Hindi through Persian. Its \textbf{\hi{ख़}}, with the dot below, is the letter Hindi writes in many words that came this way.
\end{cousinweb}

\subsection*{Your turn}
Say these aloud:

\begin{itemize}
\raggedright
  \item \emph{ākhir}, then \emph{ākhir meṁ}
  \item \textquotedblleft{}After two hours the train came at last\textquotedblright{} — \emph{do ghaṇṭe ke bād ākhir ṭren āī}
  \item \textquotedblleft{}In the end we reached home\textquotedblright{} — \emph{ākhir meṁ ham ghar pahũce}
\end{itemize}

\begin{tcolorbox}[breakable,colback=teal!4,colframe=teal!35!black,title={Before you move on}]
What does \textbf{\hi{आख़िर}} mean? (\textquotedblleft{}At last, finally\textquotedblright{}.) \textquotedblleft{}In the end\textquotedblright{}? (\textbf{\hi{आख़िर} \hi{में}}.)
\end{tcolorbox}
\section[\texorpdfstring{\hi{पहले} … \hi{फिर} … (pahle … phir …)}{पहले … फिर … (pahle … phir …)}]{\hi{पहले} … \hi{फिर} … (pahle … phir …) — first, then, after that, in the end}

\label{lesson:HI-C279-pahlephir}



\begin{quote}
You know \textbf{\hi{पहले}}, \textquotedblleft{}first, earlier\textquotedblright{}, and \textbf{\hi{फिर}}, \textquotedblleft{}again, then\textquotedblright{}. Say both. Then \textquotedblleft{}in the end\textquotedblright{}. (\emph{ākhir meṁ}.)
\end{quote}

\subsection*{You'll want to know: the order words}
\begin{quote}
\textbf{\hi{पहले} \hi{मैं} \hi{स्टेशन} \hi{गया}, \hi{फिर} \hi{मैंने} \hi{टिकट} \hi{लिया।}}
\end{quote}

\begin{quote}
\emph{pahle maiṁ sṭeśan gayā, phir maiṁne ṭikaṭ liyā.} — \textquotedblleft{}First I went to the station, then I bought a ticket.\textquotedblright{}
\end{quote}

A story marks its steps:

\noindent
\begin{tabularx}{\linewidth}{@{}>{\raggedright\arraybackslash}X>{\raggedright\arraybackslash}X@{}}
\toprule
\textbf{step} & \textbf{the word} \\
\midrule
the first & \textbf{\hi{पहले}} (\emph{pahle}) \\
the next & \textbf{\hi{फिर}} (\emph{phir}) \\
the one after & \textbf{\hi{उसके} \hi{बाद}} (\emph{uske bād}) \\
the last & \textbf{\hi{आख़िर} \hi{में}} (\emph{ākhir meṁ}) \\
\bottomrule
\end{tabularx}

Each one opens its sentence, and the verb after it is the plain past. \textbf{\hi{टिकट}} is masculine, so \textbf{\hi{लिया}}.

\subsection*{You'll want to know: \hi{उसके} \hi{बाद}}
\textbf{\hi{उसके} \hi{बाद}} is \textquotedblleft{}after that\textquotedblright{}: \textbf{\hi{उस}}, the form \textbf{\hi{वह}} takes before a postposition, as in \textbf{\hi{उसका}}, then \textbf{\hi{के} \hi{बाद}}. \textbf{\hi{इसके} \hi{बाद}} (\emph{iske bād}), \textquotedblleft{}after this\textquotedblright{}, is built the same way, on \textbf{\hi{इस}}, the form \textbf{\hi{यह}} takes before a postposition.

\begin{quote}
\textbf{\hi{उसके} \hi{बाद} \hi{हम} \hi{गाड़ी} \hi{में} \hi{बैठे।}}
\end{quote}

\begin{quote}
\emph{uske bād ham gāṛī meṁ baiṭhe.} — \textquotedblleft{}After that we got on the train.\textquotedblright{}
\end{quote}

\subsection*{Your turn}
Say these aloud:

\begin{itemize}
\raggedright
  \item \textquotedblleft{}First we ate, then we went to the market\textquotedblright{} — \emph{pahle hamne khānā khāyā, phir ham bāzār gae}
  \item \textquotedblleft{}After that I went home\textquotedblright{} — \emph{uske bād maiṁ ghar gayā}
\end{itemize}

\begin{tcolorbox}[breakable,colback=teal!4,colframe=teal!35!black,title={Before you move on}]
Which word opens the first step? (\textbf{\hi{पहले}}.) The next? (\textbf{\hi{फिर}}.) \textquotedblleft{}After that\textquotedblright{}? (\textbf{\hi{उसके} \hi{बाद}}.) The last? (\textbf{\hi{आख़िर} \hi{में}}.) Which past follows them? (The plain past.)
\end{tcolorbox}
\section[(reading)]{(reading) — a summer holiday in the village}

\label{lesson:HI-C279-chuttikadin}



\begin{quote}
Eighty-nine words, every one met before. Mira is writing. She went to her village last summer, and she tells it in order.
\end{quote}

\begin{tcolorbox}[breakable,skin=enhanced,colback=orange!6,colframe=orange!45!black,boxrule=0.5pt,arc=1mm,left=6pt,right=6pt,top=4pt,bottom=4pt,fonttitle=\bfseries,title={Read this}]
\begin{quote}
\hi{पिछले} \hi{साल} \hi{गर्मी} \hi{की} \hi{छुट्टी} \hi{में} \hi{मैं} \hi{अपने} \hi{गाँव} \hi{गई।} \hi{मुझे} \hi{गाँव} \hi{बहुत} \hi{पसंद} \hi{है।} \hi{पहले} \hi{मैंने} \hi{स्टेशन} \hi{पर} \hi{टिकट} \hi{लिया।} \hi{फिर} \hi{मैं} \hi{गाड़ी} \hi{में} \hi{बैठी।} \hi{गाड़ी} \hi{में} \hi{मैं} \hi{एक} \hi{किताब} \hi{पढ़} \hi{रही} \hi{थी।} \hi{अचानक} \hi{तेज़} \hi{बारिश} \hi{हुई।} \hi{शाम} \hi{को} \hi{मैं} \hi{गाँव} \hi{पहुँची।} \hi{वहाँ} \hi{दादी} \hi{खाना} \hi{बना} \hi{रही} \hi{थीं।} \hi{मैंने} \hi{हाथ} \hi{धोकर} \hi{दादी} \hi{के} \hi{साथ} \hi{खाना} \hi{खाया।} \hi{खाना} \hi{खाने} \hi{के} \hi{बाद} \hi{दादी} \hi{ने} \hi{मुझे} \hi{एक} \hi{कहानी} \hi{सुनाई।} \hi{उसके} \hi{बाद} \hi{मैं} \hi{छत} \hi{पर} \hi{गई।} \hi{आसमान} \hi{में} \hi{बहुत} \hi{तारे} \hi{थे।} \hi{आख़िर} \hi{में} \hi{मैं} \hi{थककर} \hi{सो} \hi{गई।} \hi{वह} \hi{दिन} \hi{बहुत} \hi{अच्छा} \hi{था।}
\end{quote}

\emph{Read it:} the passage once without stopping; do not translate

\emph{Read it:} the passage again, and sort every past verb: a scene or an event
\end{tcolorbox}

\begin{culture}
\textbf{\hi{गई}}, \textbf{\hi{बैठी}} and \textbf{\hi{पहुँची}} end in \textbf{-\hi{ई}}: Mira is a woman, and these verbs have no thing. \textbf{\hi{मैंने} \hi{टिकट} \hi{लिया}} and \textbf{\hi{मैंने} … \hi{खाना} \hi{खाया}} follow the thing, and \textbf{\hi{टिकट}} and \textbf{\hi{खाना}} are masculine. \textbf{\hi{दादी} \hi{ने} \hi{मुझे} \hi{एक} \hi{कहानी} \hi{सुनाई}} follows \textbf{\hi{कहानी}}, which is feminine; \textbf{\hi{मुझे}} is the one it was told to.

\textbf{\hi{पढ़} \hi{रही} \hi{थी}} and \textbf{\hi{बना} \hi{रही} \hi{थीं}} are the scene. \textbf{\hi{थीं}}, with the dot, is for the grandmother, said with respect. \textbf{\hi{तेज़} \hi{बारिश}} is heavy rain: \textbf{\hi{तेज़}} is \textquotedblleft{}fast, sharp\textquotedblright{}, and rain can be sharp.
\end{culture}

\subsection*{Your turn}
Say these aloud:

\begin{itemize}
\raggedright
  \item what Mira did first, at the station
  \item what her grandmother was doing when she arrived
  \item how the day ended, with \emph{ākhir meṁ}
\end{itemize}

\begin{tcolorbox}[breakable,colback=teal!4,colframe=teal!35!black,title={Before you move on}]
Why is \textbf{\hi{बना} \hi{रही} \hi{थीं}} in the scene form? (The cooking was going on when Mira arrived.) Which words keep the steps in order? (\textbf{\hi{पहले}}, \textbf{\hi{फिर}}, \textbf{\hi{उसके} \hi{बाद}}, \textbf{\hi{आख़िर} \hi{में}}.)
\end{tcolorbox}
\section[(review)]{Review — a story in order}

\label{lesson:HI-R279-a-story-in-order}



\begin{quote}
Say \textquotedblleft{}I washed my hands and ate\textquotedblright{}. (\emph{maiṁne hāth dhokar khānā khāyā}.) Then tell this morning in two steps, with \textbf{\hi{पहले}} and \textbf{\hi{फिर}}.
\end{quote}

\subsection*{Your turn}
Say these aloud:

\begin{itemize}
\raggedright
  \item First I got up, then I bathed.
  \item After breakfast I went to the station and bought a ticket.
  \item I was reading on the train when it suddenly rained.
  \item After that we sat in the courtyard, and in the end we slept.
  \item Mira's holiday in four sentences, with pahle, phir, uske bād and ākhir meṁ.
\end{itemize}

\begin{tcolorbox}[breakable,colback=teal!4,colframe=teal!35!black,title={Before you move on}]
\textquotedblleft{}I ate and fell asleep\textquotedblright{}? (\textbf{\hi{मैं} \hi{खाना} \hi{खाकर} \hi{सो} \hi{गया}}: no \textbf{\hi{ने}}, because \textbf{\hi{सो} \hi{गया}} has no thing.) \textquotedblleft{}After eating\textquotedblright{}? (\textbf{\hi{खाना} \hi{खाने} \hi{के} \hi{बाद}}.) \textquotedblleft{}After that\textquotedblright{}? (\textbf{\hi{उसके} \hi{बाद}}.) \textquotedblleft{}At last\textquotedblright{}? (\textbf{\hi{आख़िर}}.) Who told Mira a story, and why \textbf{\hi{सुनाई}}? (Her grandmother: \textbf{\hi{दादी} \hi{ने} … \hi{कहानी} \hi{सुनाई}}, and \textbf{\hi{कहानी}} is feminine.)
\end{tcolorbox}
